% !TEX root = ./adlt_homework_8.tex
\documentclass[12pt]{article}

% --- Core math + proof tools ---
\usepackage{amsmath, amssymb, amsthm}

% --- Lists and spacing ---
\usepackage{enumitem}
\setlist[enumerate,1]{label=\textbf{(\alph*)}, leftmargin=2em, itemsep=0.25em}
\setlist[itemize]{leftmargin=1.5em, itemsep=0.25em}
\setlength{\parskip}{0.35em}
\setlength{\parindent}{0pt}

% --- Page geometry + header/footer ---
\usepackage{geometry}
\geometry{margin=1in}
\usepackage{csquotes} % to use blockquotes
\usepackage{graphicx} % for figures
\graphicspath{{images/}} % tells LaTeX to look in the “images” folder

\usepackage{fancyhdr}
\pagestyle{fancy}
\fancyhf{}
\rhead{MATH 421 — 003, Fall 2025}
\lhead{Homework 8}
\cfoot{\thepage}
\renewcommand{\headrulewidth}{0.4pt}

% --- Links (nice for references, stays subtle) ---
\usepackage[hidelinks]{hyperref}

% --- Title info ---
\title{Homework 8}
\author{Alejandro De La Torre \\ MATH 421 — The Theory of Single Variable Calculus \\ Section 003 \\ Instructor: Grace Work}
\date{November 2025}

% --- (Optional) proof environments you might want ---
\newtheorem*{claim}{Claim}
\newtheorem*{lemma}{Lemma}
\newtheorem*{remark}{Remark}

% --- Convenience: a Problem header command ---
\newcommand{\problem}[2][]{%
  \vspace{0.5em}
  \hrule\vspace{0.4em}
  \textbf{Problem #2}\if\relax\detokenize{#1}\relax\else\; \textit{(#1)}\fi
  \vspace{0.25em}\hrule\vspace{0.8em}
}

% --- Convenience: an ignore command to comment out blocks ---
\newcommand{\ignore}[1]{}


\begin{document}
\maketitle

% ------------------ collaboration + sources ------------------
\section*{Collaboration Statement}
% (List any classmates you collaborated with here, or write ``I worked alone'' if you did not collaborate with anyone.)

\bigskip


% ------------------ problem 1 ------------------
\problem[Cauchy Sequence: $s_n = \frac{1}{n}$]{1}
Prove, directly from the definitions, that ($s_n$) where $s_n = \frac{1}{n}$ is a Cauchy sequence.
\vspace{0.75em}

\textbf{Discussion.} \\
So I initially wanted to prove that $s_n = \frac{1}{n}$ is a converges and then say by 10.9 Lemma that since it converges, it must be Cauchy. However, the problem specifically asks to prove it directly from the definitions, so I will need to work with the definition of a Cauchy sequence directly. So to review, here is the definition of a Cauchy sequences: 

\vspace{0.25em}

\textbf{10.8 Definition.}\\
\textit{A sequence ($s_n$) of real numbers is called a \textit{Cauchy sequence} if for each $\varepsilon > 0$, there exists an $N$ such that $m, n >N$ implies $|s_n - s_m| < \varepsilon$.}

\vspace{0.25em}

My plan is to show that for any given $\varepsilon > 0$, I can find an $N$ such that for all $m, n \geq N$, the absolute difference $|s_n - s_m|$ is less than $\varepsilon$. I will manipulate the expression $|s_n - s_m| = \left|\frac{1}{n} - \frac{1}{m}\right|$ and find a suitable $N$ based on $\varepsilon$ that guarantees that $\frac{1}{n} + \frac{1}{m} < \varepsilon$ for all $m, n \geq N$.

To find such an $N$, notice that for $m, n \geq N$,
\[
\left| \frac{1}{n} - \frac{1}{m} \right|
= \left| \frac{1}{n} + (\frac{-1}{m}) \right|
\leq \frac{1}{n} + \frac{1}{m}.
\]
By triangle inequality.

Thus, if we can make $\frac{1}{m} + \frac{1}{n} < \varepsilon$, the Cauchy condition will hold. So the question becomes: how can we choose an $N$ such that for all $m, n \geq N$, we guarantee $\frac{1}{m} + \frac{1}{n} < \varepsilon$? Well, we would just need to ensure that both $\frac{1}{m} < \frac{\varepsilon}{2}$ and $\frac{1}{n} < \frac{\varepsilon}{2}$. This would guarantee that their sum is less than $\varepsilon$.

Then we would just manipulate these inequalities to find a suitable $N$. More concretely though, for $\frac{1}{n} < \frac{\varepsilon}{2}$, how big would $n$ need to be for $\frac{1}{n} < \frac{\varepsilon}{2}$ to hold? Rearranging gives $n > \frac{2}{\varepsilon}$. Similarly, for $\frac{1}{m} < \frac{\varepsilon}{2}$, we get $m > \frac{2}{\varepsilon}$. Thus, if we choose $N$ to be any natural number greater than $\frac{2}{\varepsilon}$, then for all $m, n \geq N$, both inequalities will hold.

But then how could we know that we can find such an $N$? We know we can find a natural number $N$ such that $N > \frac{2}{\varepsilon}$? Well, by the Archimedean property, for any real number $\frac{2}{\varepsilon}$ there exists a natural number $N$ such that 
\[
N > \frac{2}{\varepsilon}.
\]
Then, for all $m, n \geq N$ we have $\frac{1}{m} \leq \frac{\varepsilon}{2}$ and $\frac{1}{n} \leq \frac{\varepsilon}{2}$, so that
\[
\frac{1}{n} + \frac{1}{m} \leq \frac{\varepsilon}{2} + \frac{\varepsilon}{2} = \varepsilon.
\]
Therefore, this choice of $N$ guarantees that $|s_n - s_m| < \varepsilon$ for all $m, n \geq N$.


\vspace{0.25em}

\textbf{Proof.}\\
Let $\varepsilon >0$ and $N > \frac{2}{\varepsilon}$. Then, for all $m, n > N$,
\[
|s_n - s_m| = \left| \frac{1}{n} - \frac{1}{m} \right| = \left| \frac{1}{n} + \left( \frac{-1}{m} \right) \right| \leq \frac{1}{n} + \frac{1}{m} < \frac{\varepsilon}{2} + \frac{\varepsilon}{2} = \varepsilon.
\]

Since $\varepsilon > 0$ was arbitrary, ($s_n$) is a Cauchy sequence.
\qed

% ------------------ problem 2 ------------------
\problem[Cauchy Sequence: $x_n = \sqrt{n}$]{2}
Let $x_n = \sqrt{n}$. Show that while the difference between consecutve terms goes to 0, i.e., show $\lim |x_{n+1} - x_n| = 0$, the sequence ($x_n$) is not a Cauchy sequence.
\vspace{0.75em}

\textbf{Discussion.} \\
%(Very brief plan / key ideas.)
\vspace{0.25em}

\textbf{Proof.}\\
% (Your proof goes here.)
\qed

% ------------------ problem 3 ------------------
\problem[Ross Chapter 2 10.6 (a)]{3}
Let ($s_n$) be a sequence such that $|s_{n+1} - s_n| < 2^{-n}$ for all $n \in \mathbb{N}$. Prove that ($s_n$) is a Cauchy sequence and hence a cin
\vspace{0.75em}

\textbf{Discussion.} \\
%(Very brief plan / key ideas.)
\vspace{0.25em}

\textbf{Proof.}\\
% (Your proof goes here.)
\qed

% ------------------ problem 4 ------------------
\problem[Ross Chapter 2 14.6]{4}
\begin{enumerate}[label=(\alph*)]
  \item Prove that if $\sum |a_n|$ converges and ($b_n$) is a bounded sequence, then $\sum a_n b_n$ converges. \textit{Hint: Use Theorem 14.4.}
  \item Observe that Corollary 14.7 is a special case of part (a).
\end{enumerate}
\vspace{0.75em}

\textbf{Discussion.} \\
%(Very brief plan / key ideas.)
\vspace{0.25em}

\textbf{Proof.}\\
% (Your proof goes here.)
\qed

% ------------------ problem 5 ------------------
\problem[Ross Chapter 2 14.7]{5}
Prove that if $\sum a_n$ is a convergent series of nonnegative numbers and $p > 1$, then $\sum a_{n}^p$ converges.
\vspace{0.75em}

\textbf{Discussion.} \\
%(Very brief plan / key ideas.)
\vspace{0.25em}

\textbf{Proof.}\\
% (Your proof goes here.)
\qed

% ------------------ problem 6 ------------------
\problem[$\varepsilon - \delta$ Convgernce Proof]{6}
Prove using the $\varepsilon - \delta$ definition of the limit of a function that if $\lim_{x \to a} f(x) = L$, then $\lim_{x \to a} 4f(x) = 4L$.
\vspace{0.75em}

\textbf{Discussion.} \\
%(Very brief plan / key ideas.)
\vspace{0.25em}

\textbf{Proof.}\\
% (Your proof goes here.)
\qed


\end{document}